%% LyX 2.4.4 created this file.  For more info, see https://www.lyx.org/.
%% Do not edit unless you really know what you are doing.
\documentclass[english,t]{beamer}
\usepackage{mathptmx}
\usepackage[T1]{fontenc}
\usepackage[latin9]{inputenc}
\usepackage{units}
\usepackage{amssymb}
\usepackage{cancel}

\makeatletter
%%%%%%%%%%%%%%%%%%%%%%%%%%%%%% Textclass specific LaTeX commands.
% this default might be overridden by plain title style
\newcommand\makebeamertitle{\frame{\maketitle}}%
% (ERT) argument for the TOC
\AtBeginDocument{%
  \let\origtableofcontents=\tableofcontents
  \def\tableofcontents{\@ifnextchar[{\origtableofcontents}{\gobbletableofcontents}}
  \def\gobbletableofcontents#1{\origtableofcontents}
}

%%%%%%%%%%%%%%%%%%%%%%%%%%%%%% User specified LaTeX commands.
\usetheme{Warsaw}
% or ...

\setbeamercovered{transparent}
% or whatever (possibly just delete it)
\setbeamercovered{invisible} %makes overlays invisible


%gets rid of navigation symbols
\setbeamertemplate{navigation symbols}{}

\usepackage{multicol}

\usepackage{tikz}
\usetikzlibrary{calc}
\usetikzlibrary{plotmarks}
\usepackage{pgfplots}
\pgfplotsset{compat=newest}

\usepackage{calc}
\usepackage[nomessages]{fp}

\usepackage{advdate}    % Advancing/saving dates
\usepackage{datetime}   % Dates formatting
\usepackage{datenumber} % Counters for dates
\usepackage[super]{nth}

%\usepackage{pgfpages}
%\pgfpagesuselayout{4 on 1}[a4paper, landscape, border shrink=7mm]
%\pgfpageslogicalpageoptions{1}{border code=\pgfusepath{stroke}}
%\pgfpageslogicalpageoptions{2}{border code=\pgfusepath{stroke}}
%\pgfpageslogicalpageoptions{3}{border code=\pgfusepath{stroke}}
%\pgfpageslogicalpageoptions{4}{border code=\pgfusepath{stroke}}
%%%Don't forget to add 'handout'

\makeatother

\usepackage{babel}
\begin{document}
\newcommand{\xmax}{2000}
\newcommand{\xmin}{0}
\newcommand{\xstep}{0} %must be xmin+n
\newcommand{\ymax}{500}
\newcommand{\ymin}{0}
\newcommand{\ystep}{0} %must be ymin+n

\newcounter{countEx} \setcounter{countEx}{0}
\newcounter{countProb} \setcounter{countProb}{0}
\resetcounteronoverlays{countEx} \resetcounteronoverlays{countProb}

\newcommand{\ex}[3]{
	\addtocounter{countEx}{1}
	\only<#1-#2>{\begin{exampleblock} {Example \chNum.\secNum.\arabic{countEx}.} #3	\end{exampleblock}}
}

\newcommand{\prob}[4]{
	\addtocounter{countProb}{1}
	\only<#1-#2>{\begin{exampleblock} {Problem  \chNum.\secNum.\arabic{countProb}.\,#3} #4	\end{exampleblock}}
}

\newcommand{\yline}[4]{
	(#2 - #4)/(#1 - #3)*(x - #1) + #2
}

\beamerdefaultoverlayspecification{<+->}

%%%%Chapter and Section numbers%%%%%
\newcommand{\chNum}{5}
\newcommand{\secNum}{2 and 5.3}
\newcommand{\secTitle}{Future and Present Value of an Annuity}

\title[Section \chNum.\secNum: \secTitle]{Section \chNum.\secNum: \secTitle}

\subtitle{Finite Mathematics~~MTH 132~~Spring 2014}

\author{Dr.\,James Ashe}

\titlegraphic{\includegraphics[height=1.5cm]{/home/jim/JamesAshe/JCSU/images/jcsu_bull}}

\institute{Johnson C. Smith University}

\date{{}}

\makebeamertitle 
\begin{frame}{5.2 Future Value of an Annuity}

\begin{definition}
{}An \emph{annuity }is an any continuing payment with a fixed total
annual amount. 

\begin{itemize}
\item [ex.]<2->Home mortgage, car payments, regular deposits to a retirement
account, etc.
\end{itemize}
\end{definition}

\begin{columns}


\column{5.5cm}

\ex{3}{}{If you deposit \$1,000 into a savings account at the end of every year at a rate of 8\% compounded yearly, how much money will be in the account after 5 years?}

\column{.5cm}

\vspace{1cm}

\column{5cm}
\begin{itemize}
\item [yr 1:]<4->$1000(1.08)^{4}=1360.49$
\item [yr 2:]<5->$1000(1.08)^{3}=1259.71$
\item [yr 3:]<6->$1000(1.08)^{2}=1166.40$
\item [yr 4:]<7->$1000(1.08)^{1}=1080.00$
\item [yr 5:]<8->$1000=1000.00$
\item [TOTAL:]<9-> ${\displaystyle \sum_{n=0}^{5}100(1.08)^{n-1}=}5866.60$
\end{itemize}
\end{columns}
\end{frame}

\begin{frame}{5.2 Future Value of an Annuity}

\begin{itemize}
\item []$S=$future value
\item []$R=$periodic payment
\item []$i=$interest rate per period $(i=\nicefrac{r}{m})$
\item []$n=$number of periods \vspace{-.5cm}
\end{itemize}
\begin{columns}


\column{5.5cm}
\begin{block}<2->{Future Value of an Annuity}
$S=R\left[\frac{\left(1+i\right)^{n}-1}{i}\right]$
\end{block}
\vfill \ex{3}{}{If you deposit \$1,000 into a savings account at the end of every year at a rate of 8\% compounded yearly, how much money will be in the account after 5 years?}

\column{5.5cm}
\begin{block}<4->{Sinking Fund Payment}
$R=\frac{S\cdot i}{(1+i)^{n}-1}$
\end{block}
\ex{5}{}{You need to save \$15,000 in 10 years. What size monthly payments should you make into an account which pays \%8 compounded monthly? (payments at end of month)}
\end{columns}
\end{frame}

\begin{frame}{5.3 Present Value of an Annuity; Amortization}

\begin{columns}


\column{5.5cm}

\ex{3}{}{Your college fund pays \$5,000 every every for 4 years. If the funds are in an account which accrues \%5 compounded yearly, how much was originally in the account?}

\column{.5cm}

\vspace{1cm}

\column{5cm}
\begin{itemize}
\item [yr 1:]<2->$17,729.75(1.05)-5000=13616.23$
\item [yr 2:]<3->$13616.23(1.05)-5000=9297.05$
\item [yr 3:]<4->$9297.05(1.05)-5000=4761.90$
\item [yr 4:]<5->$4761.90(1.05)-5000=0$
\end{itemize}
\end{columns}
\end{frame}

\begin{frame}{5.3 Present Value of an Annuity; Amortization}

\begin{columns}


\column{5.5cm}

\ex{1}{}{You buy a house for \$220,000 with a 30 year mortgage at rate of \%6 compounded monthly with payments of \$1,319.01.}

\column{.5cm}

\vspace{1cm}

\column{5cm}
\begin{itemize}
\item Amount owed:
\item [month 1:]<2->$219,780.9$
\item [month 2:]<3->$219,560.7$
\item []<4->$\vdots$
\item [month 12:]<4->$217,298.18$ (after \$15,828.12 in payments)
\item []<6->$\vdots$
\item [month 358:]<7->$2,618.37$
\item [month 359:]<8->$1,312,45$
\end{itemize}
\end{columns}
\end{frame}

\begin{frame}{5.3 Present Value of an Annuity; Amortization}

\begin{itemize}
\item []$P=$present value of annuity
\item []$R=$periodic payment
\item []$i=$interest rate per period $(i=\nicefrac{r}{m})$
\item []$n=$number of periods $\left(n=mt\right)$\vspace{-.5cm}
\end{itemize}
\begin{columns}


\column{5.5cm}
\begin{block}<2->{Amortization Payments}
$R=\frac{P}{\left[1-(1+i)^{-n}\right]}$
\end{block}
\vfill \ex{3}{}{You buy a house for \$220,000 with a 30 year mortgage at rate of \%6 compounded monthly. What are the payments?}

\column{5.5cm}
\begin{block}<4->{Present Value of an Annuity}
$P=R\left[\frac{1-\left(1+i\right)^{-n}}{i}\right]$
\end{block}
\ex{5}{}{Your college fund pays \$5,000 every every 6  months for 4 years. If the funds are in an account which accrues \%5 compounded semiannually, how much was originally in the account?}
\end{columns}
\end{frame}

\end{document}
